\documentclass[11pt]{article}
\usepackage[a4paper,margin=18mm]{geometry}
\usepackage{fontspec}
\setmainfont{Latin Modern Roman}
\usepackage{amsmath,amssymb,amsthm,amscd,graphicx,color}
\usepackage{xurl}
\usepackage[unicode,hidelinks]{hyperref}
\allowdisplaybreaks
\setlength\emergencystretch{3em}

\numberwithin{equation}{section}
\newcommand{\underaccent}[2]{\mathord{\underset{\smash{\scriptscriptstyle\sim}}{#2}}}


\makeatletter
\newtheorem{theorem}{Theorem}
\newtheorem*{theorem*}{Theorem}
\newtheorem{lemma}{Lemma}
\newtheorem*{lemma*}{Lemma}
\newtheorem{corollary}{Corollary}
\newtheorem*{corollary*}{Corollary}
\newtheorem{proposition}{Proposition}
\newtheorem*{proposition*}{Proposition}
\newtheorem{conjecture}{Conjecture}
\newtheorem*{conjecture*}{Conjecture}
{\theoremstyle{definition} \newtheorem{definition}{Definition}
\newtheorem*{definition*}{Definition}
\newtheorem{example}{Example}
\newtheorem*{example*}{Example}
\newtheorem{remark}{Remark}
\newtheorem*{remark*}{Remark}
\newtheorem{note}{Note}
\newtheorem*{note*}{Note}
}
\makeatother

\makeatletter
\def\bra#1{\left\langle #1\right|}
\def\ket#1{\left| #1\right\rangle}
\def\bracket#1#2{\left\langle #1  \,\right| \left. #2 \right\rangle}
\def\FL#1{\left\lfloor #1 \right\rfloor}
\makeatother
\begin{document}
\begin{center}\Large Irreducible successive singular-vector quotients for integer-spin d=\allowbreak{}1 complex conformal Galilei algebras at arbitrary D-weight and nonzero P\_ell weight\end{center}
\noindent\textbf{Stable ID:} 2000\par
\noindent\textbf{Authors:} Naruhiko Aizawa; Radhakrishnan Chandrashekar; Jambulingam Segar\par
\noindent\textbf{Original work:} Lowest Weight Representations, Singular Vectors and Invariant Equations for a Class of Conformal Galilei Algebras\par
\noindent\textbf{Version:} Official SIGMA 11 (2015) article 002, DOI 10.3842/SIGMA.2015.002; journal PDF and native TeX acquired 2026-10-01, exact bytes pinned by SHA256\par
\noindent\textbf{Selected task scope:} Only the nonzero-p branch of Theorem1: complex Lie algebra g\_ell with d=\allowbreak{}1 and positive integer ell, arbitrary lowest D-weight delta, nonzero P\_ell-weight p. The spin1 Verma module and the explicit successive quotient for ell>=\allowbreak{}2 are irreducible; the latter has basis H\textasciicircum{}k P\_(ell-1)\textasciicircum{}m u0 for k,m>=\allowbreak{}0. No zero-p conclusion selected.\par
\noindent\textbf{Difficulty:} H2: The author proof constructs a chain of subsingular-vector quotients for every integer spin and computes the graded Shapovalov determinants in those quotients. A final argument excludes singular vectors at every remaining level, giving an infinite-dimensional irreducible module with an explicit two-generator PBW basis.\par
\noindent\textbf{Editorial boundary note (not author text):} Only the complete original nonzero-p branch of Theorem 1 is selected; the zero-p branch is excluded without alteration, keeping the original enclosing theorem/list boundaries. All algebra/Verma/Shapovalov definitions and the full spin-one determinant proof are retained. For every integer spin ell>=\allowbreak{}2, all successive singular/subsingular quotient constructions, exact quotient PBW bases, graded determinant induction, unique next singular-vector calculation and final exclusion at every remaining level are supplied together as one irreducible-module outcome. Standard PBW and lowest-weight/submodule principles remain background. Original matrix-entry and superscript notation are preserved. No zero-p/sl2 conclusion or invariant-PDE family is selected.\par
\noindent\textbf{Source:} \url{https://sigma-journal.com/2015/002/}\quad DOI: \url{https://doi.org/10.3842/SIGMA.2015.002}\par
\noindent\textbf{License and attribution:} Copyright retained by the named authors. Published by SIGMA under CC BY 4.0: \url{https://creativecommons.org/licenses/by/4.0/}. Source: \url{https://sigma-journal.com/about.html}. Adaptation consists of standalone excerpt formatting, cover and numbering restoration; original author language and mathematical text are retained.\par
\noindent\textbf{Editorial symbol notice (not author text):} The text calls each interim submodule the largest submodule while subsequently taking further quotients. The selected module is fixed by the explicit successive construction and final Lemma5 relations, not by treating any interim quotient as already irreducible. Original pages are unchanged.\par
\noindent\textbf{Typesetting adaptation (not author text):} The unavailable accents package is replaced by a standard-AMS under-tilde typesetting macro; every source use has tilde as its accent argument. Every original vector/index invocation is retained unchanged.\par
\medskip\hrule\medskip\noindent\textbf{Author text begins below.}\par
\setcounter{section}{1}
\setcounter{subsection}{0}
\setcounter{subsubsection}{0}
\setcounter{equation}{0}
% Original source lines 130--459
\section[$d=1$ Conformal Galilei algebras ${\mathfrak g}_{\ell}$]{$\boldsymbol{d=1}$ Conformal Galilei algebras $\boldsymbol{{\mathfrak g}_{\ell}}$}

The complex Lie algebra $ {\mathfrak g}_{\ell} $ for a~f\/ixed integer $ \ell $ has the elements~\cite{NdORM1}:
\begin{gather*}
D, \ H, \ C, \ P_n, \quad n = 0, 1,   \dots,   2\ell.
%\label{generators}
\end{gather*}
Their nonvanishing commutators are given~by
\begin{alignat}{4}
& [D, H] = H, \qquad & & [D, C] = - C, \qquad & & [C, H] = 2D,& \nonumber\\
& [H, P_n] = -n P_{n-1}, \qquad & & [D, P_n] = (\ell-n) P_n, \qquad & & [C, P_n] = (2\ell-n) P_{n+1}.&\label{nonvanishing-commutators}
\end{alignat}
One may see from this that $ \langle   P_0, P_1, \dots, P_{2\ell}   \rangle $ is an Abelian ideal of $ {\mathfrak
g}_{\ell} $, so that $ {\mathfrak g}_{\ell} $ is \textit{not} semisimple.
It is known that this Lie algebra has no central extensions~\cite{MT}.
The subalgebra spanned by $ \langle
 H, D, C
\rangle $ is isomorphic to $ {\mathfrak{so}}(2,1) \simeq {\mathfrak{sl}}(2,{\mathbb R}) \simeq {\mathfrak{su}}(1,1)$.
The Abelian subalgebra spanned by $ \langle
P_n
\rangle_{n=0, 1, \dots, 2\ell} $ carries the spin $ \ell $ representation of the $ {\mathfrak{sl}}(2,{\mathbb R})$
subalgebra.

This algebra may be realized as generators of transformation of $(1+1)$-dimensional spacetime
\begin{gather*}
H = \frac{\partial}{\partial t},
\qquad
D =-t\frac{\partial}{\partial t} -\ell x \frac{\partial}{\partial x},
\qquad
C = t^2 \frac{\partial}{\partial t} + 2 \ell tx \frac{\partial}{\partial x},
\qquad
P_n = (-t)^n \frac{\partial}{\partial x}.
%\label{CGArealization-in-space}
\end{gather*}
In this realization,~$H$,~$D$ and~$C$ generates time translation, dilatation and the special conformal representation,
respectively.
Meanwhile the generator of space translation is represented by~$P_0$, the Galilei transformation generator by~$P_1$,
the transformation to a~reference frame with constant acceleration is given by~$P_2$ and so on.

One may introduce the algebraic anti-involution $ \omega: {\mathfrak g}_{\ell} \to {\mathfrak g}_{\ell} $~by
\begin{gather*}
\omega(D) = D, \qquad \omega(H) = C, \qquad \omega(P_n) = P_{2\ell-n}.
%\label{anti-inv}
\end{gather*}
It is not dif\/f\/icult to verify that $ \omega $ satisf\/ies the required relations:
\begin{gather*}
\omega([X,Y]) = [\omega(Y), \omega(X)], \qquad \omega^2(X) = X \qquad  \forall\, X \in {\mathfrak g}_{\ell}.
%\label{anti-inv2}
\end{gather*}
Let us def\/ine the degree of the generators based on their commutator with respect to~$D$:
\begin{gather*}
\deg(D) = 0, \qquad \deg(H) = 1, \qquad \deg(C) = -1, \qquad \deg(P_n) = \ell-n.
%\label{degree}
\end{gather*}
With respect to the sign of the degree one may def\/ine the triangular decomposition of ${\mathfrak g}_{\ell}$:
\begin{gather*}
    {\mathfrak g}_{\ell} = {\mathfrak g}_{\ell}^+ \oplus {\mathfrak g}_{\ell}^0 \oplus {\mathfrak g}_{\ell}^-,
%\label{tri-decomp}
\end{gather*}
where
\begin{gather*}
{\mathfrak g}_{\ell}^+ = \langle H, P_0, P_1, \dots, P_{\ell-1} \rangle,
\qquad
{\mathfrak g}_{\ell}^0 = \langle D, P_{\ell} \rangle,
\qquad
{\mathfrak g}_{\ell}^- = \langle C, P_{\ell+1}, P_{\ell+2}, \dots, P_{2\ell} \rangle.
\end{gather*}
This is a~decomposition of $ {\mathfrak g}_{\ell} $ as a~direct sum of the vector spaces.


\section{Verma modules and Shapovalov determinant}
\label{SEC:VM}

We start with the one-dimensional module $ {\mathbb C} \ket{\delta,p} $ over $ {\mathfrak b} = {\mathfrak g}_{\ell}^0
\oplus {\mathfrak g}_{\ell}^-$. The lowest weight vector $ \ket{\delta,p} $ is def\/ined~by
\begin{gather*}
D \ket{\delta,p} = \delta \ket{\delta,p},
\qquad
P_{\ell} \ket{\delta,p} = p \ket{\delta,p},
\qquad
X \ket{\delta,p} = 0,
\qquad
 \forall \, X \in {\mathfrak g}_{\ell}^-.
%\label{LWV}
\end{gather*}
Then for each pair of $(\delta,p) $ the Verma module over $ {\mathfrak g}_{\ell} $ is def\/ined as usual (see, e.g.,~\cite{Dix}): $ V^{\delta,p}_{\ell} = U({\mathfrak g}_{\ell}) \otimes_{U({\mathfrak b})} \ket{\delta,p} $ where $U({\mathfrak g}_{\ell})$ and $ U({\mathfrak b}) $ are the enveloping algebras of $ {\mathfrak g}_{\ell} $ and $
\mathfrak b$, respectively.
In order to specify the basis of $ V^{\delta,p}_{\ell} $ we introduce the $ \ell$-component vector
\begin{gather*}
\underaccent{\tilde}{m} = (m_1, m_2, \dots, m_{\ell})  \in {\mathbb R}^{\ell},
%\label{m-vec-def}
\end{gather*}
where $ m_i$ $(i=1,2,\dots,\ell) $ are non-negative integers.
With this the basis of $ V^{\delta,p}_{\ell} $ is given~by
\begin{gather}
\ket{k,\underaccent{\tilde}{m}} = H^k P_{\ell-1}^{m_1} P_{\ell-2}^{m_2} \cdots P_0^{m_{\ell}} \ket{\delta,p}.
\label{anyl-basis}
\end{gather}
We also introduce the~$\ell$-component vectors
\begin{gather*}
\underaccent{\tilde}{\epsilon}_j = (0, \dots, 0, 1, 0, \dots,0) \in {\mathbb R}^{\ell},
\qquad
j=1,2,\dots, \ell.
%\label{tilde-epsilon-def}
\end{gather*}
where the~$j$th entry of $\underaccent{\tilde}{\epsilon}_j $ is~1, and all other entries are~0.
It may not be dif\/f\/icult to prove the following relation by induction on $k$:
\begin{gather}
[P_n, H^k] = \sum\limits_{i=1}^{\min\{n,k\}} \binom{n}{i} \frac{k!}{(k-i)!}  H^{k-i}   P_{n-i},
\qquad
n \geq 1.
\label{PnHkcomm}
\end{gather}
It follows that the action of $ {\mathfrak g}_{\ell}^0$ and ${\mathfrak g}_{\ell}^+$ on $\ket{k,\underaccent{\tilde}{m}}$:
\begin{gather*}
    D \ket{k,\underaccent{\tilde}{m}} = \left(\delta + k+ \sum\limits_{i=1}^{\ell} i   m_i\right) \ket{k,\underaccent{\tilde}{m}},
\nonumber
\\
    P_{\ell} \ket{k,\underaccent{\tilde}{m}} = p \ket{k,\underaccent{\tilde}{m}} + \sum\limits_{i=1}^{\min\{\ell,k\}}
\binom{\ell}{i} \frac{k!}{(k-i)!} \ket{k-i,\underaccent{\tilde}{m}+\underaccent{\tilde}{\epsilon}_{i}},
\nonumber
\\
    H \ket{k,\underaccent{\tilde}{m}}= \ket{k+1,\underaccent{\tilde}{m}},
\nonumber
\\
    P_n \ket{k,\underaccent{\tilde}{m}} = \sum\limits_{i=1}^{\min\{n,k\}} \binom{n}{i} \frac{k!}{(k-i)!}
\ket{k-i,\underaccent{\tilde}{m}+\underaccent{\tilde}{\epsilon}_{\ell-n+i}},
\qquad
1 \leq n \leq \ell-1,
\nonumber
\\
    P_0 \ket{k,\underaccent{\tilde}{m}} = \ket{k,\underaccent{\tilde}{m}+\underaccent{\tilde}{\epsilon}_{\ell}}.
%\label{g0p-action-any}
\end{gather*}
From this one see that $ V^{\delta,p}_{\ell} $ has a~grading structure according to the eigenvalue of~$D$:
\begin{gather*}
    V^{\delta,p}_{\ell} = \bigoplus_{N=0}^{\infty} \big(V^{\delta,p}_{\ell}\big)_N,
\qquad
\big(V^{\delta,p}_{\ell}\big)_N = \big\{ \ket{v} \in V^{\delta,p}_{\ell} \, | \, D \ket{v} = (\delta+N) \ket{v} \big\}.
%\label{Verma-grading}
\\
    N = k+ \sum\limits_{i=1}^{\ell} i   m_i.
%\label{N-definition}
\end{gather*}
We refer to~$N$ as the \textit{level} as usual.

Now we def\/ine Shapovalov form on a~module over $ {\mathfrak g}_{\ell} $~\cite{Shapo} (see also, e.g.,~\cite{KacRai}).
That is a~contravariant Hermitian form $ \bracket{\cdot}{\cdot} $ def\/ined by making use of the anti-involution~ $\omega$.
For the Verma module $ V^{\delta,p}_{\ell} $ it is def\/ined as follows:   Let $ \ket{x} $ and $ \ket{y} $ be any two
vectors in $ V^{\delta,p}_{\ell}$.
They may be written as
\begin{gather*}
\ket{x} = X \ket{\delta,p},
\qquad
\ket{y} = Y \ket{\delta,p},
\qquad
X, Y \in U({\mathfrak g}_{\ell}^+).
\end{gather*}
Then
\begin{gather}
\bracket{x}{y} = \bra{\delta,p} \omega(X) Y \ket{\delta,p},
\qquad
\bracket{\delta,p}{\delta,p} = 1.
\label{innerprod-def}
\end{gather}
Next we def\/ine the Shapovalov determinant at level~$N$.
Let $ \ket{v_1}, \ket{v_2}, \dots, \ket{v_r} $ be a~set of basis of the subspace $ \big(V^{\delta,p}_{\ell}\big)_N$.
We consider the matrix $ (\bracket{v_i}{v_j}) $ whose entries are     the Shapovalov form and we call the determinant of
this matrix, the Shapovalov determinant at level $N: \Delta^{(\ell)}_N = \det (\bracket{v_i}{v_j})$.
We give an explicit formulae of the Shapovalov determinants of $ {\mathfrak g}_{\ell}$, since $ \Delta^{(\ell)}_N $ is
important to know the reducibility of $V^{\delta,p}_{\ell}$.
\begin{proposition}\label{prop1}
The Shapovalov determinants $ \Delta^{(\ell)}_N$ at level~$N$ of ${\mathfrak g}_{\ell}$ are given as follows $($up to
overall sign$)$:
\begin{alignat}{3}
   & i)
\quad &&
\Delta_N^{(1)} = \left(\prod\limits_{m=0}^N m! \right)^2 (2p)^{N(N+1)},&
\label{KDl1}
\\
  &   ii)
\quad &&
\Delta^{(\ell)}_N = (\ell+1)^2   p^2   \delta_{N1},
\qquad
\ell \geq 2.&
%\label{KDell}
\end{alignat}
\end{proposition}

\begin{proof}
The case with $ \ell = 1 $ shows a~deviation from other values of $ \ell$.
We treat the case with $ \ell = 1 $ separately.
The basis of $ V^{\delta,p}_1 $ is specif\/ied by two nonnegative integers:
\begin{gather*}
\ket{k,m} = H^k P_0^m \ket{\delta,p}.
\end{gather*}
The next relation obtained from~\eqref{PnHkcomm} is useful in the following computation:
\begin{gather}
P_2 \ket{k,m} = 2kp \ket{k-1,m} + k(k-1) \ket{k-2,m+1}.
\label{ActionP2}
\end{gather}
The level is $ N = k +m $ so that the basis of $ \big(V^{\delta,p}_1\big)_N $ is given~by
\begin{gather*}
\ket{m,N-m},
\qquad
0 \leq m \leq N.
\end{gather*}
The product of two vectors in $ \big(V^{\delta,p}_1\big)_N $ is
\begin{gather}
\bracket{N-k,k}{m,N-m} = \bra{N-k,0} P_0^{N-m} P_2^k \ket{m,0}.
\label{IPl=1}
\end{gather}

It follows from~\eqref{ActionP2} that $P_2^k \ket{m,0}$ is a~linear combination of $ \ket{m-k-j,j} $ with $0 \leq j
\leq k$.
Therefore, $ P_2^k \ket{m,0} = 0$ if $k > m$.
Thus we have proved the following lemma:

\begin{lemma}
\label{lemma1}
$ \bracket{N-k,k}{m,N-m} = 0 $ if $ k > m$.
\end{lemma}

By def\/inition $ \Delta^{(1)}_N $ is given by (up to sign)
\begin{gather*}
\Delta^{(1)}_N = \left|
\begin{matrix}
 \bracket{N,0}{0,N} & \bracket{N,0}{1,N-1} & \dots & \bracket{N,0}{N,0}
\\
\bracket{N-1,0}{0,N} & \bracket{N-1,0}{1,N-1} & \dots & \bracket{N-1,0}{N,0}
\\
\vdots & \vdots & \ddots & \vdots
\\
\bracket{0,N}{0,N} & \bracket{0,N}{1,N-1} & \dots & \bracket{0N}{N,0}
\end{matrix}
\right|.
\end{gather*}
By Lemma~\ref{lemma1}, $ \Delta^{(1)}_N $ is the determinant of upper triangular matrix.
Thus
\begin{gather*}
\Delta_N^{(1)} = \prod\limits_{m=0}^N \bracket{N-m,m}{m,N-m}.
\end{gather*}
Each factor is calculated by~\eqref{IPl=1} as follows
\begin{gather*}
\bracket{N-m,m}{m,N-m} = \bra{N-m,m} P_0^{N-m} P_2^m \ket{N-m,m} = (N-m)!   m!   (2p)^N,
\end{gather*}
where we used the relation obtained from~\eqref{ActionP2}: $ P_2^m \ket{m,0} = (2p)^m m ! \ket{0,0}$.
This completes the proof of~\eqref{KDl1}.

Now let us turn to the case with $\ell \geq 2$.
In this case $ \Delta^{(\ell)}_N \neq 0 $ only if $N = 1$.
As we shall see, this fact stems from that at least two rows of $ \Delta^{(\ell)}_N $ are proportional if $ N > 1$.
The basis of $(V^{\delta,p}_{\ell})_N $ is given by~\eqref{anyl-basis} with the constraint $  N = k+
\sum\limits_{i=1}^{\ell} i   m_i$.
First we prove the following lemma:
\begin{lemma}
\label{lemma2}
In the subspace $ (V^{\delta,p}_{\ell})_N $ we have the relation
\begin{gather*}
\bracket{0,\underaccent{\tilde}{n}}{k,\underaccent{\tilde}{m}} = f^{(\ell)}_{N,\underaccent{\tilde}{n}}(p) \delta_{kN},
%\label{IPlge2}
\end{gather*}
where $ f^{(\ell)}_{N,\underaccent{\tilde}{n}}(p) $ is a~function of~$p$ determined by $ \ell$, $N $ and $
\ket{0,\underaccent{\tilde}{n}}$.
\end{lemma}
\begin{proof}
By def\/inition
\begin{gather}
    \bracket{0,\underaccent{\tilde}{n}}{k,\underaccent{\tilde}{m}} = \bra{\delta,p} P_{2\ell}^{n_{\ell}}
P_{2\ell-1}^{n_{\ell-1}} \cdots P_{\ell+1}^{n_1} H^k P_{\ell-1}^{m_1} \cdots P_0^{m_{\ell}} \ket{\delta,p},
\label{IPlge2-2}
\\
    N = \sum\limits_{i=1}^{\ell} i n_i = k + \sum\limits_{i=1}^{\ell} i  m_i.
\nonumber
\end{gather}
Because of~\eqref{PnHkcomm} the right hand side of~\eqref{IPlge2-2} will be a~linear combination of the terms
\begin{gather*}
\bra{\delta,p} H^{k'} P_{\ell-1}^{m'_1} \cdots P_0^{m'_{\ell}} \ket{\delta,p},
\end{gather*}
with the condition $ m'_k \geq m_k $ for all~$k$.
However, such terms give nonvanishing contributions only if
\begin{gather*}
k' = m'_1 = m'_2 = \dots = m'_{\ell} = 0.
\end{gather*}
This implies $ m_k = 0 $ for all $k$, i.e., $ k = N$.
\end{proof}

It follows from Lemma~\ref{lemma2} that two rows in $ \Delta^{(\ell)}_N $ labelled by $ \bra{0,\underaccent{\tilde}{n}}
$ and $ \bra{0,\underaccent{\tilde}{n'}} $ are proportional.
If $ N \geq 2 $ then there exists at least one such pair of rows, but no such pair for $N = 1$.
Thus $ \Delta^{(\ell)}_N = 0 $ if $ N \geq 2$.
On the other hand the $ N = 1 $ subspace is two-dimensional with the basis $ \ket{1, \underaccent{\tilde}{0}}, \ket{0,
\underaccent{\tilde}{\epsilon}_1} $ where $\underaccent{\tilde}{0}$ denotes the zero vector in ${\mathbb R}^{\ell}$.
Thus $ \Delta^{(\ell)}_1 $ is computed as follows
\begin{gather*}
\Delta^{(\ell)}_1 = \left|
\begin{matrix} \bracket{1, \underaccent{\tilde}{0}}{1, \underaccent{\tilde}{0}} & \bracket{1, \underaccent{\tilde}{0}}{0,
\underaccent{\tilde}{\epsilon}_1}
\vspace{1mm}\\  \bracket{0, \underaccent{\tilde}{\epsilon}_1}{1, \underaccent{\tilde}{0}} & \bracket{0,
\underaccent{\tilde}{\epsilon}_1}{0, \underaccent{\tilde}{\epsilon}_1}
\end{matrix}
\right| = \left|
\begin{matrix} 2 \delta & (\ell+1)p
\vspace{1mm}\\  (\ell+1)p & 0
\end{matrix}
\right| = - (\ell+1)^2   p^2   .
\end{gather*}
This completes the proof of Proposition~\ref{prop1}.
\end{proof}



\begin{proposition}\qquad
\label{prop2}
\begin{enumerate}\itemsep=0pt
\item[$i)$] $ V^{\delta,p}_{\ell} $ is irreducible if $ \ell = 1 $ and $ p \neq 0$.
\item[$ii)$] $ V^{\delta,p}_{\ell} $ is reducible if $ \ell \geq 2 $ and $ p \neq 0$.
\item[$iii)$] $ V^{\delta,0}_{\ell} $ is reducible for all values of $ \ell$.
\end{enumerate}
\end{proposition}



\begin{proof}
(i) is the corollary of Proposition~\ref{prop1}(i).
Proposition~\ref{prop1} also suggests the existence of singular vectors in $V^{\delta,0}_{\ell} $ for any $\ell$ and
in $V^{\delta,p}_{\ell}$ with $\ell \geq 2$ and $p \neq 0$.
In the next section we shall show that this is indeed the case.
Thus we establish~(ii) and~(iii).
\end{proof}
\setcounter{section}{4}
\setcounter{subsection}{0}
\setcounter{subsubsection}{0}
\setcounter{equation}{0}
% Original source lines 467--491
Before giving the list let us recall the def\/inition of singular vectors.
A~singular vector $\ket{v_s} \in V^{\delta,p}_{\ell}$ is yet another lowest weight vector which is not proportional to
$\ket{\delta,p}$.
Namely, $ \ket{v_s} $ satisf\/ies the conditions:
\begin{gather}
    \ket{v_s} \neq {\mathbb C} \ket{\delta,p},
\qquad\!\!
    D \ket{v_s} = \delta' \ket{v_s},
\qquad\!\!
P_{\ell} \ket{v_s} = p' \ket{v_s},
\qquad\!\!
X \ket{v_s} = 0,
\qquad\!\!
\forall \,  X \in {\mathfrak g}_{\ell}^-.\!\!\!
\label{SV-def}
\end{gather}
According to Proposition~\ref{prop1}, singular vectors in $V^{\delta,p}_{\ell}$ may be dif\/ferent for $p \neq 0$ and
$p = 0$.
We treat these cases separately.

\subsection[$p \neq 0$]{$\boldsymbol{p \neq 0}$}

In this case the Verma modules over $ \ell = 1 $ algebra have no singular vectors.
For the algebra $ {\mathfrak g}_{\ell} $ with $ \ell \geq 2 $, singular vectors may exist only in the subspace $
\big(V^{\delta,p}_{\ell}\big)_N $ with $ N \geq 2$.
\setcounter{section}{5}
\setcounter{subsection}{2}
\setcounter{subsubsection}{0}
\setcounter{equation}{5}
% Original source lines 886--907
\section[Irreducible lowest weight modules of ${\mathfrak g}_{\ell}$]{Irreducible lowest weight modules of $\boldsymbol{{\mathfrak g}_{\ell}}$}
\label{SEC:IRREP}

We have shown that the Verma modules over $ {\mathfrak g}_{\ell} $ are reducible in many cases
(Proposition~\ref{prop2}).
It is known that the Verma module is, in a~sense, the largest lowest weight module.
That is, one can derive all irreducible lowest weight modules starting from the Verma module $ V^{\delta,p}_{\ell}$.
The purpose of this section is to obtain some types of irreducible lowest weight modules explicitly.
Our results are summarized in the next theorem.

\begin{theorem}
\label{theorem2}
The lowest weight modules over ${\mathfrak g}_{\ell}$ given below are irreducible:
\begin{itemize}\itemsep=0pt
\item $ p \neq 0 $
\begin{enumerate}\itemsep=0pt
\item[$i)$] the Verma module $ V^{\delta,p}_1 $ for $ \ell = 1$;
\item[$ii)$] the quotient module $ V_{\ell}^{(\ell)} $ for $ \ell \geq 2$.
This module is infinite-dimensional with the basis vectors $ H^k P_{\ell-1}^m \big|u_0^{(\ell)}\big\rangle $, where $ k$, $m $ are
nonnegative integers.
See Lemma~{\rm \ref{lemma5}} for the definition of the lowest weight vector $ \big| u_0^{(\ell)}\big\rangle$.
\end{enumerate}
\setcounter{section}{6}
\setcounter{subsection}{0}
\setcounter{subsubsection}{0}
\setcounter{equation}{0}
% Original source lines 915--1286
\end{itemize}
\end{theorem}

Theorem~\ref{theorem2} coincides with the results in~\cite{LuMazoZhao}.

\subsection[Proof for $p \neq 0$]{Proof for $\boldsymbol{p \neq 0}$}
\label{PFqn0}

The Verma modules $ V^{\delta,p}_{\ell} $ are reducible for $ \ell \geq 2 $ so we restrict ourselves to $ \ell \geq 2$.
We consider the quotient module $ V^{\delta,p}_{\ell}/{\cal I}^{(2)} $ where $ {\cal I}^{(2)} $ is the largest
${\mathfrak g}_{\ell}$-submodule of $ V^{\delta,p}_{\ell}$.
Since there is no singular vectors in the $ N = 1 $ subspace $ (V^{\delta,p}_{\ell})_1$, $ {\cal I}^{(2)} $ will be
induced by the singular vector in the $ N = 2 $ subspace.
\begin{lemma}
\label{lemma3}
There exists precisely one $($up to an overall constant$)$ singular vector in the $ N = 2 $ subspace $
(V^{\delta,p}_{\ell})_2$.
This singular vector is given~by
\begin{gather*}
\big|v_s^{(2)}\big\rangle = (2p (\ell+1) P_{\ell-2} - (\ell+2) P_{\ell-1}^2) \ket{\delta,p}.
%\label{SVN2}
\end{gather*}
\end{lemma}


\begin{proof}
The basis of $ \big(V^{\delta,p}_{\ell}\big)_2 $ is $ \ket{2,\underaccent{\tilde}{0}}, \ket{1,\underaccent{\tilde}{\epsilon}_1},
\ket{0,2\underaccent{\tilde}{\epsilon}_1}, \ket{0,\underaccent{\tilde}{\epsilon}_2}$.
The singular vector $ \big|v_s^{(2)}\big\rangle$ is a~linear combination of the basis:
\begin{gather*}
\big|v_s^{(2)}\big\rangle  = c_1 \ket{2,\underaccent{\tilde}{0}} + c_2 \ket{1,\underaccent{\tilde}{\epsilon}_1} + c_3
\ket{0,2\underaccent{\tilde}{\epsilon}_1} + c_4 \ket{0,\underaccent{\tilde}{\epsilon}_2}.
\end{gather*}
It must satisfy the condition
\begin{gather*}
0 = P_{\ell+1} \big|v_s^{(2)}\big\rangle  = 2 (\ell+1)  p   c_1 \ket{1,\underaccent{\tilde}{0}} + (\ell+1)   (\ell  c_1 + p
c_2) \ket{0,\underaccent{\tilde}{\epsilon}_1}.
\end{gather*}
Thus $ c_1 = c_2 = 0$.
Furthermore, the condition $ C \big|v_s^{(2)}\big\rangle  = 0 $ yields the relation
\begin{gather*}
2 (\ell+1)   p   c_3 + (\ell+2) c_4 = 0.
\end{gather*}
This proves Lemma~\ref{lemma3}.
\end{proof}

Def\/ine $ {\cal I}^{(2)} = U({\mathfrak g}_{\ell}^{+}) \big|v_s^{(2)}\big\rangle $, then $ {\cal I}^{(2)} $ is the largest $
{\mathfrak g}_{\ell}$-submodule in $ V^{\delta,p}_{\ell}$.
Let $ \big|u_0^{(2)}\big\rangle  $ be the lowest weight vector in $ V^{\delta,p}_{\ell}/{\cal I}^{(2)}$.
Then
\begin{gather*}
    D \big|u_0^{(2)}\big\rangle  = \delta \big|u_0^{(2)}\big\rangle ,
\qquad
P_{\ell} \big|u_0^{(2)}\big\rangle  = p \big|u_0^{(2)}\big\rangle ,
\nonumber
\\
    X \big|u_0^{(2)}\big\rangle  = 0, \qquad X \in {\mathfrak g}_{\ell}^-,
\qquad
P_{\ell-2} \big|u_0^{(2)}\big\rangle  = \frac{\ell+2}{2p (\ell+1)} P_{\ell-1}^2 \big|u_0^{(2)}\big\rangle .
%\label{LWVQuo2}
\end{gather*}
It follows that the basis of $ V^{\delta,p}_{\ell}/{\cal I}^{(2)} $ is given~by
\begin{gather*}
\big|k,\underaccent{\tilde}{m}^{(2)}\big\rangle  = H^k   P_{\ell-1}^{m_1}  P_{\ell-3}^{m_3} \cdots P_{0}^{m_{\ell}}
\big|u_0^{(2)}\big\rangle ,
%\label{BasisQuo2}
\end{gather*}
where $ \underaccent{\tilde}{m}^{(2)} = (m_1, 0, m_3, \dots,m_{\ell}) \in {\mathbb R}^{\ell}$.
Observing the relation
\begin{gather*}
D \big|k,\underaccent{\tilde}{m}^{(2)}\big\rangle  = \left(\delta+k +m_1 + \sum\limits_{i=3}^{\ell}i m_i\right)
\big|k,\underaccent{\tilde}{m}^{(2)}\big\rangle ,
\end{gather*}
we def\/ine the level $ N^{(2)} $ in the quotient space $ V^{\delta,p}_{\ell}/{\cal I}^{(2)} $ by $  N^{(2)}
= k + m_1 + \sum\limits_{i=3}^{\ell} i m_{i}$.
The vectors $ \{
\ket{1,\underaccent{\tilde}{0}},
\ket{0,\underaccent{\tilde}{\epsilon}_1}
\} $ and $ \{
\ket{2,\underaccent{\tilde}{0}},
\ket{1,\underaccent{\tilde}{\epsilon}_1},
\ket{0,2\underaccent{\tilde}{\epsilon}_1}
\}$ form a~basis of $ N^{(2)} = 1 $ and $ N^{(2)} = 2 $ subspaces of $ V^{\delta,p}_{\ell}/{\cal I}^{(2)}$,
respectively.
The Shapovalov determinant for $ N^{(2)} = 1 $ subspace is same as~\eqref{SEC:VM} and given by $ (\ell+1)^2 p^2$.
While that for $ N^{(2)} = 2 $ subspace is calculated as follows
\begin{gather*}
    \left|
\begin{matrix} \bracket{2,\underaccent{\tilde}{0}}{2,\underaccent{\tilde}{0}} &
\bracket{2,\underaccent{\tilde}{0}}{1,\underaccent{\tilde}{\epsilon}_1} &
\bracket{2,\underaccent{\tilde}{0}}{0,2\underaccent{\tilde}{\epsilon}_1}
\vspace{1mm}\\ \bracket{1,\underaccent{\tilde}{\epsilon}_1}{2,\underaccent{\tilde}{0}} &
\bracket{1,\underaccent{\tilde}{\epsilon}_1}{1,\underaccent{\tilde}{\epsilon}_1} &
\bracket{1,\underaccent{\tilde}{\epsilon}_1}{}
\vspace{1mm}\\ \bracket{0,2\underaccent{\tilde}{\epsilon}_1}{2,\underaccent{\tilde}{0}} &
\bracket{0,2\underaccent{\tilde}{\epsilon}_1}{1,\underaccent{\tilde}{\epsilon}_1} &
\bracket{0,2\underaccent{\tilde}{\epsilon}_1}{0,2\underaccent{\tilde}{\epsilon}_1}
\end{matrix}
\right|
\\
\qquad{}
= \left|
\begin{matrix} 4(2\delta+1) \delta & 2(2\delta+1)(\ell+1)p & 2(\ell+1)^2 p^2
\vspace{1mm}\\ 2(2\delta+1)(\ell+1)p & (\ell+1)^2 p^2 & 0
\vspace{1mm}\\ 2(\ell+1)^2 p^2 & 0 & 0
\end{matrix}
\right| = -4(\ell+1)^6 p^6.
\end{gather*}
Therefore there exist no singular vectors in $ N^{(2)} = 1, 2 $ subspaces.
On the other hand, one f\/inds a~singular vector in the $ N^{(2)} = 3 $ subspace.

\begin{lemma}
\label{lemma4}
There exists precisely one $($up to overall constant$)$ singular vector in the $ N^{(2)} = 3 $ subspace of $
V^{\delta,p}_{\ell}/{\cal I}^{(2)}$.
The singular vector is given~by
\begin{gather*}
\big|v_s^{(3)}\big\rangle  = \big(3! p^2 (\ell+1)^2  P_{\ell-3} - (\ell+2)  (\ell+3)  P_{\ell-1}^3\big) \big|u_0^{(2)}\big\rangle .
%\label{SVN3}
\end{gather*}
\end{lemma}

\begin{proof}
The lemma can be proved in a~way exactly similar to Lemma~\ref{lemma3}.
The basis of the level $ N^{(2)} = 3 $ subspace is given~by
\begin{gather*}
\ket{3,\underaccent{\tilde}{0}},
\qquad
\ket{2,\underaccent{\tilde}{\epsilon}_1},
\qquad
\ket{1,2\underaccent{\tilde}{\epsilon}_1},
\qquad
\ket{0,3\underaccent{\tilde}{\epsilon}_1},
\qquad
\ket{0,\underaccent{\tilde}{\epsilon}_3}.
\end{gather*}
The singular vector $ \big|v_s^{(3)}\big\rangle  $ is a~linear combination of the basis:
\begin{gather*}
\big|v_s^{(3)}\big\rangle  = c_1 \ket{3,\underaccent{\tilde}{0}} + c_2 \ket{2,\underaccent{\tilde}{\epsilon}_1} + c_3
\ket{1,2\underaccent{\tilde}{\epsilon}_1} + c_4 \ket{0,3\underaccent{\tilde}{\epsilon}_1} + c_5
\ket{0,\underaccent{\tilde}{\epsilon}_3}.
\end{gather*}
It must satisfy the condition
\begin{gather*}
0  =  P_{\ell+1} \big|v_s^{(3)}\big\rangle
\\
\hphantom{0}{}
 =  3 p c_1 \ket{2,\underaccent{\tilde}{0}} + (3\ell c_1 + 2p c_2) \ket{1,\underaccent{\tilde}{\epsilon}_1} +
\left(\frac{(\ell+2) \ell (\ell-1)}{2(\ell+1)p} c_1 + \ell c_2 + (\ell+1) p c_3 \right) \ket{0,
2\underaccent{\tilde}{\epsilon}_1}.
\end{gather*}
Thus $ c_1 = c_2 = c_3 = 0$.
Furthermore, the condition $ C \big|v_s^{(3)}\big\rangle  = 0 $ yields the relation
\begin{gather*}
3 (\ell+1) p c_4 + \frac{(\ell+3) (\ell+2)}{2(\ell+1)p} c_5 = 0.
\end{gather*}
This proves  Lemma~\ref{lemma4}.
\end{proof}

The vector $ \big|v_s^{(3)}\big\rangle  $ is singular only in the quotient space $ V^{\delta,p}_{\ell}/{\cal I}^{(2)} $ and not in
$ V^{\delta,p}_{\ell} $ itself.
Such vector is called \textit{subsingular}~\cite{Dob95,Dob97}.

It follows from Lemma~\ref{lemma4} that the subspace $ {\cal I}^{(3)} = U({\mathfrak g}_{\ell}^+) \big|v_s^{(3)}\big\rangle  $ is
the largest ${\mathfrak g}_{\ell}$-submodule in $ V^{\delta,p}_{\ell}/{\cal I}^{(2)}$.
Now we consider the quotient space $ V^{\delta,p}_{\ell}/{\cal I}^{(2)}/{\cal I}^{(3)}:= \big(V^{\delta,p}_{\ell}/{\cal
I}^{(2)}\big)/{\cal I}^{(3)}$.
The lowest weight vector $ \big|u_0^{(3)}\big\rangle  $ of this quotient space is def\/ined~by
\begin{gather*}
    D \big|u_0^{(3)}\big\rangle  = \delta \big|u_0^{(3)}\big\rangle ,
\qquad
P_{\ell} \big|u_0^{(3)}\big\rangle  = p \big|u_0^{(3)}\big\rangle ,
\qquad
    X \big|u_0^{(3)}\big\rangle  = 0, \qquad X \in {\mathfrak g}_{\ell}^-,
\nonumber
\\
    P_{\ell-2} \big|u_0^{(3)}\big\rangle  = \frac{\ell+2}{2p  (\ell+1)} P_{\ell-1}^2 \big|u_0^{(3)}\big\rangle ,
\qquad
P_{\ell-3} \big|u_0^{(3)}\big\rangle  = \frac{(\ell+2) (\ell+3)}{3!  p^2  (\ell+1)^2} P_{\ell-1}^3\big|u_0^{(3)}\big\rangle .
%\label{LWVQuo3}
\end{gather*}
The basis of $ V^{\delta,p}_{\ell}/{\cal I}^{(2)}/{\cal I}^{(3)} $ is given~by
\begin{gather*}
H^k  P_{\ell-1}^{m_1} P_{\ell-4}^{m_4} \cdots P_0^{m_{\ell}} \big|u_0^{(3)}\big\rangle ,
\end{gather*}
and we def\/ine the level $  N^{(3)} = k + m_1 + \sum\limits_{i=4}^{\ell} i  m_i$.
With this setting one can show the followings:
\begin{enumerate}\itemsep=0pt
\item[i)] the Shapovalov determinant does not vanish for $ 1 \leq N^{(3)} \leq 3 $ so that there exist no singular vectors
in the subspaces with $ N^{(3)} = 1, 2, 3$;
\item[ii)] there exists a~unique singular vector in the level $ N^{(3)} = 4 $ subspace.
\end{enumerate}
By the singular vector in $ N^{(3)} = 4 $ subspace, one can def\/ine the largest ${\mathfrak g}_{\ell}$-submodule $ {\cal
I}^{(4)} $ in $ V^{\delta,p}_{\ell}/{\cal I}^{(2)}/{\cal I}^{(3)} $ and consider the quotient space $
V^{\delta,p}_{\ell}/{\cal I}^{(2)}/{\cal I}^{(3)}/{\cal I}^{(4)}$.

In fact, one can repeat this process until we arrive at $ V^{\delta,p}_{\ell}/{\cal I}^{(2)}/\cdots/{\cal I}^{(\ell)}$.
This is assured by the next lemma:
\begin{lemma}
\label{lemma5}
Suppose that we have arrived at the quotient space $ V_{\ell}^{(\lambda)}:= V^{\delta,p}_{\ell}/{\cal
I}^{(2)}/\cdots/{\cal I}^{(\lambda)}$, $2 \leq \lambda \leq \ell$.
Namely, we have the quotient space with the basis
\begin{gather*}
H^k P_{\ell-1}^{m_1} P_{\ell-\lambda-1}^{m_{\lambda+1}} \cdots P_0^{m_{\ell}} \big|u_0^{(\lambda)}\big\rangle ,
%\label{LWVQuoL}
\end{gather*}
where $ \big|u_0^{(\lambda)}\big\rangle  $ is the lowest weight vector in $ V_{\ell}^{(\lambda)} $, defined~by
\begin{gather}
    D \big|u_0^{(\lambda)}\big\rangle  = \delta \big|u_0^{(\lambda)}\big\rangle ,
\qquad
P_{\ell} \big|u_0^{(\lambda)}\big\rangle  = p \big|u_0^{(\lambda)}\big\rangle ,
\qquad
X \big|u_0^{(\lambda)}\big\rangle  = 0, \qquad X \in {\mathfrak g}_{\ell}^-,
\nonumber
\\      P_{\ell-a} \big|u_0^{(\lambda)}\big\rangle  = \frac{(\ell+2) (\ell+3) \cdots (\ell+a)}{a!   p^{a-1}   (\ell+1)^{a-1}}
P_{\ell-1}^a \big|u_0^{(\lambda)}\big\rangle ,
\qquad
a = 2, 3, \dots, \lambda.
\label{Assumption-lambda}
\end{gather}
Then
\begin{enumerate}\itemsep=0pt
\item[$i)$] $ V_{\ell}^{(\lambda)} $ is the graded vector space:
\begin{gather*}
    V_{\ell}^{(\lambda)} = \bigoplus_{N^{(\lambda)}=0}^{\infty} (V_{\ell}^{(\lambda)})_{N^{(\lambda)}},
\qquad
\big(V_{\ell}^{(\lambda)}\big)_{N^{(\lambda)}} = \big\{ \ket{v} \in V_{\ell}^{(\lambda)} \, | \, D \ket{v} = (\delta +
N^{(\lambda)}) \ket{v} \big\},
\nonumber
\\
    N^{(\lambda)} = k + m_1 + \sum\limits_{i=\lambda+1}^{\ell} i m_i.
%\label{gradedquotientspace}
\end{gather*}
\item[$ii)$] The subspace $ \big(V_{\ell}^{(\lambda)}\big)_{N^{(\lambda)}} $ has a~nonvanishing Shapovalov determinant if $ 1 \leq
N^{(\lambda)} \leq \lambda$.
The Shapovalov determinant is given by $($up to a~sign factor$)$
\begin{gather}
\Delta_{N^{(\lambda)}} = (p  (\ell+1))^{N^{(\lambda)}(N^{(\lambda)}+1)} \prod\limits_{k=0}^{N^{(\lambda)}} k!
\big(N^{(\lambda)}-k\big)!.
\label{KDinQuotient}
\end{gather}
This implies that there exists no singular vectors in the level $ N^{(\lambda)} $ subspaces if $ 1 {\leq} N^{(\lambda)}
{\leq} \lambda$.
\item[$iii)$] If $ 2 \leq \lambda \leq \ell-1$, then there exists precisely one $($up to overall constant$)$ singular vector in the
$ N^{(\lambda)} = \lambda + 1 $ subspace of $ V^{(\lambda)}_{\ell}$.
The singular vector is given~by
\begin{gather}
\big|v_s^{(\lambda{+}1)}\big\rangle  = \big(  (\lambda+1)!  p^{\lambda} (\ell+1)^{\lambda} P_{\ell{-}\lambda{-}1} \!- (\ell+2) (\ell+3)
\cdots (\ell+\lambda+1) P_{\ell{-}1}^{\lambda{+}1}  \big)\! \big|u_0^{(\lambda)}\big\rangle .\!\!\!\!
\label{SVinQuotientSpace}
\end{gather}
\end{enumerate}
\end{lemma}

\begin{proof}
(i) Can be trivially proved.
 (ii) Suppose that $ 1 \leq N^{(\lambda)} \leq \lambda$, then the basis of $
\big(V_{\ell}^{(\lambda)}\big)_{N^{(\lambda)}} $ is given by $ \ket{k} = H^{N^{(\lambda)}-k} P_{\ell-1}^k \big|u_0^{(\lambda)}\big\rangle
$ with $ k = 0, 1, \dots, N^{(\lambda)}$.
The equality in the next equation is up to a~sign factor:
\begin{gather}
\Delta_{N^{(\lambda)}} = \left|
\begin{matrix} \bracket{0}{N^{(\lambda)}} & \bracket{0}{N^{(\lambda)}-1} & \dots & \bracket{0}{0}
\\ \bracket{1}{N^{(\lambda)}} & \bracket{1}{N^{(\lambda)}-1} & \dots & \bracket{1}{0}
\\
\vdots & \vdots & & \vdots
\\
\bracket{N^{(\lambda)}}{N^{(\lambda)}} & \bracket{N^{(\lambda)}}{N^{(\lambda)}-1} & \vdots & \bracket{N^{(\lambda)}}{0}
\end{matrix}
\right|.
\label{KDQuoSpa}
\end{gather}
The lower triangular entry of~\eqref{KDQuoSpa} is given~by
\begin{gather}
\big\langle m \big|  N^{(\lambda)}-k\big\rangle  = \big\langle u_0^{(\lambda)}\big| P_{\ell+1}^m C^{N^{(\lambda)}-m} H^k P_{\ell-1}^{N^{(\lambda)}-k}
\big|u_0^{(\lambda)}\big\rangle ,
\qquad
k < m.
\label{Prod-m-Nk}
\end{gather}
Using the commutation relations
\begin{gather*}
\big[C, H^k\big] = 2k H^{k-1} D + k(k-1) H^{k-1},
\qquad
\big[C, P_{\ell-1}^k\big] = k (\ell+1) P_{\ell-1}^{k-1} P_{\ell},
\end{gather*}
we see that~\eqref{Prod-m-Nk} is a~linear combination of
\begin{gather}
\big\langle u_0^{(\lambda)}\big| P_{\ell+1}^m H^j P_{\ell-1}^{m-j} \big|u_0^{(\lambda)}\big\rangle ,
\qquad
\max\{k-N+m, 0\} \leq j \leq k.
\label{Prod-m-Nk-2}
\end{gather}
The lower bound for~$j$ corresponds to the two possibilities $ N-m \leq k $ or $ N-m > k$.
By the relation
\begin{gather*}
\big[H, P_{\ell+1}^m\big] = -m(\ell+1) P_{\ell+1}^{m-1} P_{\ell},
\end{gather*}
the equation~\eqref{Prod-m-Nk-2} is further reduced to
\begin{gather*}
\big\langle u_0^{(\lambda)}\big| P_{\ell+1}^{m-j} P_{\ell-1}^{m-j} \big|u_0^{(\lambda)}\big\rangle  = 0.
\end{gather*}
The last equality stems from $ m -j > 0 $ which is due to the range of~$j$ and $ k < m$.
Thus all the lower triangular entries of~\eqref{KDQuoSpa} vanish.

The diagonal entries of~\eqref{KDQuoSpa} correspond to~\eqref{Prod-m-Nk} with $ m = k$.
Thus they are expanded into a~linear combination of~\eqref{Prod-m-Nk-2}.
Because of $ m = k $ there exist precisely one term in the expansion which gives the nonvanishing contribution.
Writing the coef\/f\/icient explicitly, that term is given~by
\begin{gather*}
(N^{(\lambda)}-k) ((\ell+1) p)^{N^{(\lambda)}-k} \big\langle u_0^{(\lambda)}\big| P_{\ell+1}^k H^k \big|u_0^{(\lambda)}\big\rangle  = k!
(N^{(\lambda)}-k)!   (p  (\ell+1))^{N^{(\lambda)}}.
\end{gather*}
The formula~\eqref{KDinQuotient} immediately follows from these results.

 (iii) The basis vectors of level $ N^{(\lambda)} = \lambda + 1 $ subspace are $ P_{\ell-\lambda-1}
\big|u_0^{(\lambda)}\big\rangle  $ and $ \ket{k} =$ \linebreak  $H^{\lambda+1-k}   P_{\ell-1}^k \big|u_0^{(\lambda)}\big\rangle$ with $
k=0,1,\dots,\lambda+1$.
The singular vector is a~linear combination of these vectors:
\begin{gather*}
\big|v_s^{(\lambda+1)}\big\rangle  = \sum\limits_{k=0}^{\lambda+1} \alpha_k \ket{k} + \beta P_{\ell-\lambda-1}
\big|u_0^{(\lambda)}\big\rangle .
\end{gather*}
The condition $ P_{\ell+1} \big|v_s^{(\lambda+1)}\big\rangle  = 0 $ yields
\begin{gather*}
    (\lambda+1-n) \left(  p (\ell+1)  \alpha_n + \binom{\ell+1}{2} (\lambda+2-n) \alpha_{n-1} \right) H^{\lambda-n}
P_{\ell-1}^n \big|u_0^{(\lambda)}\big\rangle
\\
\qquad{}
+ \sum\limits_{j=2}^n \binom{\ell+1}{j+1} \frac{(\lambda+1-n+j)!}{(\lambda-n)!}   \alpha_{n-j} H^{\lambda-n} P_{\ell-j}
P_{\ell-1}^{n-j} \big|u_0^{(\lambda)}\big\rangle  = 0,
\qquad
0 \leq n \leq \lambda.
\end{gather*}
It follows that $ \alpha_j = 0 $ for $ 0 \leq j \leq \lambda$.
Thus the singular vector yields
\begin{gather*}
\big|v_s^{(\lambda+1)}\big\rangle  = \big(  \alpha_{\lambda+1} P_{\ell-1}^{\lambda+1} + \beta P_{\ell-\lambda-1}  \big)
\big|u_0^{(\lambda)}\big\rangle .
\end{gather*}
The condition $ C \big|v_s^{(\lambda+1)}\big\rangle  = 0 $ gives the relation
\begin{gather*}
(\lambda+1)!  p^{\lambda} (\ell+1)^{\lambda} \alpha_{\lambda+1} + (\ell+2) (\ell+3) \cdots (\ell+\lambda+1) \beta = 0.
\end{gather*}
This implies the uniqueness of the formula~\eqref{SVinQuotientSpace} of the singular vector.
We remark that if $ \lambda=\ell $ then the vector corresponds to $ P_{\ell-\lambda-1} $ does not exist.
Thus $ \alpha_{\lambda+1} = 0 $, so that no singular vectors at level $ \ell+1$.
\end{proof}

Now we consider the module $ V^{(\ell)}_{\ell} = V^{\delta,p}_{\ell}/{\cal I}^{(2)}/ {\cdots } /{\cal I}^{(\ell)}$.
The basis of this space is $ H^k P_{\ell{-}1}^m \big|u_0^{(\ell)}\big\rangle  $ where $ \big|u_0^{(\ell)}\big\rangle  $ is the lowest weight
vector def\/ined by~\eqref{Assumption-lambda} with $ \lambda = \ell$.
By Lemma~\ref{lemma5}, there exist no singular vectors in the subspaces of $ V^{(\ell)}_{\ell} $ labelled by $
N^{(\ell)} = 1, 2, \dots, \ell$.
For $ N^{(\ell)} \geq \ell+1 $ the singular vectors may be written as
\begin{gather*}
\big|v_s^{(\ell)}\big\rangle  = \sum\limits_{k=0}^{N^{(\ell)}} \alpha_k H^k P_{\ell-1}^{N^{(\ell)}-k} \big|u_0^{(\ell)}\big\rangle .
\end{gather*}

The condition $ P_{\ell+1} \big|v_s^{(\ell)}\big\rangle  = 0 $ yields
\begin{gather*}
\sum\limits_{k=1}^{N^{(\ell)}} \sum\limits_{j=1} \alpha_k \binom{\ell+1}{j} \frac{k!}{(k-j)!} H^{k-j} P_{\ell+1-j}
P_{\ell-1}^{N^{(\ell)}-k} \big|u_0^{(\ell)}\big\rangle  = 0.
\end{gather*}
From the above relation we can see that $ \alpha_k = 0$ $(k \neq 0)$.
We also see that $ \alpha_0 = 0 $ from
\begin{gather*}
0 = C \big|v_s^{(\ell)}\big\rangle  = \alpha_0 N^{(\ell)} (\ell+1) p P_{\ell-1}^{N^{(\ell)}-1} \big|u_0^{(\ell)}\big\rangle .
\end{gather*}
Thus there are no singular vectors in $ V^{(\ell)}_{\ell}$.
That is, $ V^{(\ell)}_{\ell} $ is an irreducible ${\mathfrak g}_{\ell}$-module.
\begin{thebibliography}{99}
\bibitem[15]{Dix}
Dixmier J., Enveloping algebras, \textit{North-Holland Mathematical Library}, Vol.~14, North-Holland Publishing Co.,
Amsterdam~-- New York~-- Oxford, 1977.

\bibitem[17]{Dob95}
Dobrev V.K., Subsingular vectors and conditionally invariant ({$q$}-deformed) equations,
\href{http://dx.doi.org/10.1088/0305-4470/28/24/014}{\textit{J.~Phys.~A: Math. Gen.}} \textbf{28} (1995), 7135--7155.

\bibitem[18]{Dob97}
Dobrev V.K., Kazhdan--{L}usztig polynomials, subsingular vectors and conditionally invariant ({$q$}-deformed) equations,
in Symmetries in Science,~{IX} ({B}regenz, 1996), Plenum, New York, 1997, 47--80.

\bibitem[27]{KacRai}
Kac V.G., Raina A.K., Bombay lectures on highest weight representations of inf\/inite-dimensional {L}ie algebras,
\textit{Advanced Series in Mathematical Physics}, Vol.~2, World Scientif\/ic Publishing Co., Inc., Teaneck, NJ, 1987.

\bibitem[28]{LuMazoZhao}
L{\"u} R., Mazorchuk V., Zhao K., On simple modules over conformal {G}alilei algebras,
\href{http://dx.doi.org/10.1016/j.jpaa.2014.02.012}{\textit{J.~Pure Appl. Algebra}} \textbf{218} (2014), 1885--1899, \href{http://arxiv.org/abs/1310.6284}{arXiv:1310.6284}.

\bibitem[30]{MT}
Martelli D., Tachikawa Y., Comments on {G}alilean conformal f\/ield theories and their geometric realization,
\href{http://dx.doi.org/10.1007/JHEP05(2010)091}{\textit{J.~High Energy Phys.}} \textbf{2010} (2010), no.~5, 091,
31~pages, \href{http://arxiv.org/abs/0903.5184}{arXiv:0903.5184}.

\bibitem[32]{NdORM1}
Negro J., del Olmo M.A., Rodr{\'{\i}}guez-Marco A., Nonrelativistic conformal groups,
\href{http://dx.doi.org/10.1063/1.532067}{\textit{J.~Math. Phys.}} \textbf{38} (1997), 3786--3809.

\bibitem[35]{Shapo}
Shapovalov N.N., A~certain bilinear form on the universal enveloping algebra of a~complex semisimple {L}ie algebra,
\href{http://dx.doi.org/10.1007/BF01077650}{\textit{Funct. Anal. Appl.}} \textbf{6} (1972), 307--312.

\end{thebibliography}
\end{document}
